\documentclass[a4paper,abstract,german]{scrreprt}

%
% Packages
%
\usepackage[german]{babel}
\usepackage[utf8]{inputenc}
\usepackage[T1]{fontenc}
\usepackage{lmodern}
\usepackage{amsmath}
\usepackage{bm}
\usepackage{amsmath}
\usepackage{amsfonts} 
\usepackage{amssymb}
\usepackage{amsthm}
\usepackage{graphicx}
\usepackage{enumitem}
\usepackage{tikz}
\usepackage[left=2.5cm,right=2.5cm,top=3cm,bottom=2cm,includeheadfoot]{geometry}

\setlength{\parindent}{0pt} 



%
% Eigene Macros
%
\newcommand{\N}{\mathbb{N}} 
\newcommand{\Z}{\mathbb{Z}} 
\newcommand{\Q}{\mathbb{Q}} 
\newcommand{\R}{\mathbb{R}} 
\newcommand{\C}{\mathbb{C}} 

%
% Beginn des eigentlichen Dokuments
%
\begin{document}


\noindent
{\begin{flushright}
%
% Hier die Namen einfüllen
%	
1. NAME \\
2. NAME \\
3. NAME
\end{flushright}}
\begin{center}
	{\textbf  {Mathematik für Informatiker 1}} \\
	{\textbf {Abgabe NUMMER. Übungsblatt}}
\end{center}\vspace{0.3cm}

%
% Aufzählung der Aufgaben
%
\begin{enumerate}
	
	
	\item[\textbf {H1.1}]
	%
	% Hier ist Platz für die Lösung einer Aufgabe
	%
	Formeln im Fließtext müssen mit \verb|$| oder \verb|\(| und \verb|\)|
	umschlossen werden, also zum Beispiel:
	$a^2 + b^2 = c^2$.
	Abgesetzte Formeln können mithilfe von 
	\verb|\[ Formel \]| oder mit \\
	\verb|\begin{align*} Formel \end{align*}|
	gesetzt werden, zum
	Beispiel:
	\begin{align*}
		\sum_{k=1}^{n} k = \frac{n(n+1)}{2}
	\end{align*}
	Längere abgesetzte Formeln lassen sich mit der align-Umbegung %\verb|\begin{align*} Formel \end{align*}|
	strukturieren:
	\begin{align*}
		(A \cdot B)_{ij}
		&= \sum_{k=1}^{n} A_{ik} \cdot B_{kj} \\
		&= \ldots
	\end{align*}
	Der \LaTeX -Code wäre in dem Fall:\\
	\verb|\begin{align*}|\\
	\verb|(A \cdot B)_{ij}| \\
	\verb|&= \sum_{k=1}^{n} A_{ik} \cdot B_{kj} \\| \\
	\verb|&= \ldots| \\
	\verb|\end{align*}| \\
	Falls Sie Formeln nummerieren wollen können sie dies indem Sie die Sterne in den obigen Formel weglassen also \verb|\begin{align} \label{SuperFormel} Formel mit nummer \end{align}|. Sie können dann die Formel mit \verb| \eqref{SuperFormel} | referenzieren.
	
	\item[\textbf {H1.2}]
	
	%
	% Hier ist Platz für die Lösung einer Aufgabe
	%
	Ein paar wichtige Symbole: $\in$, $\Rightarrow$, $\Leftarrow$, $\Leftrightarrow$, $\leq$, $\geq$, $\cup$, $\cap$, $\neq$, $\subseteq$, $\supseteq$ $\subsetneq$, $\Q$, $\R$, $\C$, $\lim\limits_{n \to \infty}$, $\int_{a}^{b}$.
	
	Mehr findet man z.B. unter:
	
	\verb|https://de.wikipedia.org/wiki/Hilfe:TeX|
	und in dem Dokument darüber wie man Formel in Wuecampus setzt.
	
	\item[\textbf {H1.3}]
	Weitere Lösung\\
	
	%
	% Hier ist Platz für die Lösung einer Aufgabe
	%
	AUFGABE
	
\end{enumerate}

\end{document}
